\documentclass[a4paper,12pt]{article}

\usepackage[T2A]{fontenc}
\usepackage[utf8]{inputenc}
\usepackage[russian]{babel}
\usepackage{PTSans}
\renewcommand{\familydefault}{\sfdefault}

\usepackage{amsmath,amssymb,mathtools}
\usepackage{icomma}
\usepackage[a4paper,left=19mm,right=19mm,top=18mm,bottom=19mm]{geometry}
\usepackage[nopatch=footnote]{microtype}
\usepackage{tikz}
\usetikzlibrary{calc,arrows.meta,positioning,shapes.geometric}
\usepackage{tkz-euclide}
\usepackage{tabularx,booktabs,longtable,array}
\usepackage{enumitem}
\usepackage{needspace}
\usepackage{hyperref}
\usepackage{fancyhdr}
\usepackage{setspace}
\usepackage{xcolor}

\definecolor{accent}{HTML}{287C7A}
\definecolor{darkaccent}{HTML}{184E4D}
\definecolor{textgray}{HTML}{30383A}
\definecolor{softgray}{HTML}{E7ECEC}
\definecolor{muted}{HTML}{667375}

\hypersetup{hidelinks}
\color{textgray}
\onehalfspacing
\raggedbottom
\setlength{\parindent}{0pt}
\setlength{\parskip}{0.45em}
\setlength{\headheight}{25pt}
\setlength{\emergencystretch}{2em}
\setlist[itemize]{leftmargin=1.5em,itemsep=0.25em,topsep=0.25em}
\setlist[enumerate]{leftmargin=2em,itemsep=0.75em,topsep=0.4em}
\renewcommand{\arraystretch}{1.28}

\pagestyle{fancy}
\fancyhf{}
\fancyhead[L]{\small\color{muted}Екатерина Скелли}
\fancyhead[R]{\small\color{muted}Подготовка к МЦКО}
\fancyfoot[C]{\small\color{muted}\thepage}
\renewcommand{\headrulewidth}{0.35pt}
\renewcommand{\headrule}{%
  \hbox to\headwidth{%
    \color{softgray}\leaders\hrule height \headrulewidth\hfill
  }%
}

\newcommand{\Section}[1]{%
  \Needspace{5\baselineskip}%
  \vspace{0.6em}%
  {\Large\bfseries\color{darkaccent}#1}\par
  \vspace{0.2em}%
  {\color{accent}\rule{0.14\linewidth}{1.2pt}}\par
  \vspace{0.4em}%
}

\newcommand{\Subsection}[1]{%
  \Needspace{3\baselineskip}%
  \vspace{0.45em}%
  {\large\bfseries\color{darkaccent}#1}\par
  \vspace{0.15em}%
}

\newcommand{\Key}[1]{\textbf{\color{darkaccent}#1}}
\newcommand{\degree}{^{\circ}}
\renewcommand{\today}{19 сентября 2026 г.}

\tikzset{
  flowstart/.style={
    ellipse,
    draw=accent,
    line width=0.9pt,
    minimum width=36mm,
    minimum height=10mm,
    align=center,
    font=\small
  },
  flowstep/.style={
    rounded corners=3pt,
    draw=accent,
    line width=0.9pt,
    text width=108mm,
    minimum height=12mm,
    align=center,
    font=\small,
    inner sep=5pt
  },
  flowchoice/.style={
    diamond,
    aspect=3.1,
    draw=accent,
    line width=0.9pt,
    text width=42mm,
    align=center,
    font=\small,
    inner sep=2pt
  },
  flowarrow/.style={
    -{Stealth[length=3mm]},
    line width=0.9pt,
    color=darkaccent
  }
}

\begin{document}

% ---------- Титульный лист ----------
\thispagestyle{empty}

\begin{center}
  \vspace*{25mm}

  {\color{accent}\rule{0.18\textwidth}{1.3pt}}\\[12mm]

  {\fontsize{28}{34}\selectfont
   \bfseries\color{darkaccent}Чек-лист}\\[7mm]

  {\fontsize{19}{24}\selectfont
   Геометрия треугольника,\\
   статистика и вероятность}\\[10mm]

  {\large Подготовка к МЦКО}\\[20mm]

  {\Large\bfseries Екатерина Скелли}\\[22mm]

  {\normalsize \today}\\[6mm]

  {\normalsize Лёвин Артём Александрович}\\[2mm]

  {\normalsize\color{muted}
   эксклюзивно для Екатерины Скелли}

  \vfill

  \begin{tikzpicture}[x=1cm,y=1cm]
    \draw[accent,line width=1.1pt]
      (-7.2,0)
      .. controls (-4.8,1.0) and (-2.4,-0.9) .. (0,0)
      .. controls (2.3,0.9) and (4.8,-0.65) .. (7.2,0.2);
    \fill[accent] (0,0) circle (1.3pt);
  \end{tikzpicture}

  \vspace*{8mm}
\end{center}

\clearpage

% ---------- Навигация ----------
\Section{Как работать с чек-листом}

Материал собран в логике занятия: сначала основные геометрические связи,
затем приёмы выбора решения, после них — статистика и вероятность.

\begin{itemize}
  \item
  Переведите условие на математический язык.\\
  Отметьте равные стороны и углы, биссектрисы, медианы, высоты,
  параллельность и перпендикулярность.

  \item
  Каждый найденный угол связывайте с конкретной причиной.\\
  Используйте сумму углов треугольника, смежные или вертикальные углы,
  свойства равнобедренного треугольника и параллельных прямых.

  \item
  В задачах на вероятность сначала выпишите все равновозможные исходы,
  затем выделите благоприятные.

  \item
  В тренировочном блоке записывайте ход решения. Если задача пока не
  поддаётся, поставьте рядом прочерк и вернитесь к ней после остальных.
\end{itemize}

\Subsection{Мини-проверка перед ответом}

\begin{tabularx}{\linewidth}{
  @{}
  >{\bfseries\color{darkaccent}}p{0.27\linewidth}
  X
  @{}
}
\toprule
Что проверить & Контрольный вопрос \\
\midrule
Чертёж &
Все ли данные условия отмечены и согласуются ли они с рисунком? \\

Углы &
Получается ли сумма углов каждого использованного треугольника
равной $180\degree$? \\

Длины &
Выполняются ли все три неравенства треугольника? \\

Обоснование &
Названо ли свойство или признак, который разрешает сделанный вывод? \\

Статистика &
Учтена ли частота каждого значения и верно ли найден объём выборки? \\

Вероятность &
Все ли исходы равновозможны и посчитаны ровно один раз? \\
\bottomrule
\end{tabularx}

\clearpage

\Section{1. Углы в треугольнике}

\Key{Сумма углов треугольника:}
\[
  \angle A+\angle B+\angle C=180\degree.
\]

\Key{Смежные углы} имеют общую сторону, а две другие стороны образуют прямую:
\[
  \alpha+\beta=180\degree.
\]

\Key{Вертикальные углы} равны. Для них нельзя автоматически утверждать, что
их сумма равна $180\degree$: это возможно только тогда, когда каждый из них
равен $90\degree$.

\Subsection{Внешний угол треугольника}

Пусть сторона $AC$ продолжена за точку $C$ до точки $D$. Тогда внешний угол
при вершине $C$ равен сумме двух внутренних углов, не смежных с ним:
\[
  \angle BCD=\angle CAB+\angle ABC.
\]

\begin{center}
\begin{tikzpicture}[scale=0.92]
  \tkzDefPoint(0,0){A}
  \tkzDefPoint(3.1,2.4){B}
  \tkzDefPoint(6.1,0){C}
  \tkzDefPoint(8.0,0){D}

  \tkzDrawSegments[thick,color=accent](A,B B,C A,C C,D)
  \tkzDrawPoints(A,B,C,D)

  \tkzLabelPoints[below](A,C,D)
  \tkzLabelPoints[above](B)

  \tkzMarkAngle[size=0.72,color=darkaccent](D,C,B)
  \tkzLabelAngle[pos=1.05](D,C,B){$\gamma$}
\end{tikzpicture}
\end{center}

\Key{Пример.} Если $\angle A=42\degree$, $\angle B=62\degree$, то
\[
  \angle C
  =180\degree-42\degree-62\degree
  =76\degree,
  \qquad
  \gamma
  =180\degree-76\degree
  =104\degree.
\]

То же значение получается быстрее:
\[
  \gamma=42\degree+62\degree=104\degree.
\]

\Key{Типичная ошибка.}
По внешнему виду рисунка нельзя решать, острый угол или тупой.
Опирайтесь на условие и вычисления.

\clearpage

\Section{2. Равнобедренный треугольник}
\enlargethispage{3\baselineskip}

Если $AB=BC$, то $AC$ — основание, а углы при основании равны:
\[
  \angle BAC=\angle ACB.
\]

Верно и обратное утверждение: против равных углов лежат равные стороны.

\begin{center}
\begin{tikzpicture}[scale=0.95]
  \tkzDefPoint(0,0){A}
  \tkzDefPoint(6,0){C}
  \tkzDefPoint(3,3.6){B}
  \tkzDefPoint(3,0){M}

  \tkzDrawSegments[thick,color=accent](A,B B,C C,A B,M)
  \tkzDrawPoints(A,B,C,M)

  \tkzLabelPoints[below](A,M,C)
  \tkzLabelPoints[above](B)

  \tkzMarkSegments[mark=|,color=darkaccent](A,B B,C)
  \tkzMarkSegments[mark=||,color=darkaccent](A,M M,C)
  \tkzMarkRightAngle[size=0.28,color=darkaccent](B,M,A)
\end{tikzpicture}
\end{center}

В равнобедренном треугольнике медиана $BM$, проведённая к основанию,
одновременно является
\[
  \text{медианой},
  \qquad
  \text{биссектрисой},
  \qquad
  \text{высотой}.
\]

Следовательно,
\[
  AM=MC,
  \qquad
  \angle ABM=\angle MBC,
  \qquad
  BM\perp AC.
\]

\Subsection{Как выбрать неизвестную величину}

Если биссектриса делит угол $A$ пополам, удобно положить
\[
  \angle BAD=\angle DAC=x,
\]
тогда весь угол
\[
  \angle BAC=2x.
\]

Такая запись обычно проще, чем обозначать весь угол через $x$ и постоянно
использовать $x/2$.

\Key{Пример.}
В равнобедренном $\triangle ABC$ дано $AB=BC$,
точка $D$ лежит на стороне $BC$, $AD$ — биссектриса угла $A$,
$\angle ADC=132\degree$.

Пусть
\[
  \angle BAD=\angle DAC=x.
\]

Тогда
\[
  \angle A=2x,
  \qquad
  \angle C=2x.
\]

Из $\triangle ADC$:
\[
  x+2x+132\degree=180\degree,
  \qquad
  x=16\degree.
\]

Поэтому
\[
  \angle A=\angle C=32\degree,
\]
а
\[
  \angle B
  =180\degree-32\degree-32\degree
  =116\degree.
\]

\clearpage

\Section{3. Прямоугольный треугольник}

\Subsection{Угол $30\degree$}

Катет, лежащий напротив угла $30\degree$, равен половине гипотенузы:
\[
  \angle A=30\degree
  \quad\Longrightarrow\quad
  BC=\frac{AC}{2}.
\]

Обратное утверждение также верно: если катет равен половине гипотенузы,
то против него лежит угол $30\degree$.

\begin{center}
\begin{tikzpicture}[scale=0.95]
  \tkzDefPoint(0,0){B}
  \tkzDefPoint(6,0){A}
  \tkzDefPoint(0,3.45){C}

  \tkzDrawSegments[thick,color=accent](A,B B,C C,A)
  \tkzDrawPoints(A,B,C)

  \tkzLabelPoints[below](A,B)
  \tkzLabelPoints[above](C)

  \tkzMarkRightAngle[size=0.3,color=darkaccent](A,B,C)
  \tkzMarkAngle[size=0.65,color=darkaccent](C,A,B)
  \tkzLabelAngle[pos=0.95](C,A,B){$30\degree$}
\end{tikzpicture}
\end{center}

\Key{Пример.}
Если катет напротив угла $30\degree$ равен $10$, то гипотенуза равна $20$.

\Subsection{Угол $45\degree$}

Если один острый угол прямоугольного треугольника равен $45\degree$,
второй тоже равен $45\degree$.
Такой треугольник равнобедренный, поэтому его катеты равны.

\Subsection{Теорема Пифагора}

Для катетов $a$, $b$ и гипотенузы $c$:
\[
  c^2=a^2+b^2.
\]

Гипотенуза лежит напротив прямого угла и является самой длинной стороной.

\clearpage

\Section{4. Свойства и вспомогательные построения}

\Subsection{Признаки равенства треугольников}

\begin{tabularx}{\linewidth}{
  @{}
  >{\bfseries\color{darkaccent}}p{0.20\linewidth}
  X
  @{}
}
\toprule
Признак & Достаточные данные \\
\midrule
Первый & Две стороны и угол между ними. \\
Второй & Сторона и два прилежащих к ней угла. \\
Третий & Три стороны. \\
\bottomrule
\end{tabularx}

\Key{Полезный приём.}
Если в условии нет готовых треугольников, иногда нужно провести
вспомогательный отрезок: соединить точки, опустить перпендикуляр или
продолжить сторону. После этого ищите равные углы и общую сторону.

\Subsection{Точка на биссектрисе угла}

Любая точка на биссектрисе равноудалена от сторон угла.
Расстояния до сторон угла измеряются по перпендикулярам.

\begin{center}
\begin{tikzpicture}[scale=0.9]
  \tkzDefPoint(0,0){A}
  \tkzDefPoint(6,2.4){B}
  \tkzDefPoint(6,-2.4){D}
  \tkzDefPoint(3.0,0){C}

  \tkzDefPointBy[projection=onto A--B](C)
  \tkzGetPoint{E}

  \tkzDefPointBy[projection=onto A--D](C)
  \tkzGetPoint{F}

  \tkzDrawSegments[thick,color=accent](A,B A,D A,C C,E C,F)
  \tkzDrawPoints(A,C,E,F)

  \tkzLabelPoints[left](A)
  \tkzLabelPoints[above](C,E)
  \tkzLabelPoints[below](F)

  \tkzMarkRightAngle[size=0.25,color=darkaccent](C,E,A)
  \tkzMarkRightAngle[size=0.25,color=darkaccent](A,F,C)

  \draw[darkaccent,line width=0.8pt]
    (A) ++(0:0.55)
    arc[start angle=0,end angle=21.8,radius=0.55];

  \draw[darkaccent,line width=0.8pt]
    (A) ++(-21.8:0.55)
    arc[start angle=-21.8,end angle=0,radius=0.55];
\end{tikzpicture}
\end{center}

Здесь
\[
  CE=CF.
\]

\Subsection{Касательная и радиус}

Касательная к окружности перпендикулярна радиусу, проведённому в точку
касания. Угол между ними равен
\[
  90\degree.
\]

\clearpage

\Subsection{Неравенство треугольника}

Для сторон $a$, $b$, $c$ должны одновременно выполняться условия
\[
  a+b>c,
  \qquad
  a+c>b,
  \qquad
  b+c>a.
\]

Если одна сторона неизвестна, удобно использовать равносильную двойную оценку:
\[
  |a-b|<c<a+b.
\]

Например, для сторон $3$, $5$, $c$ получаем
\[
  2<c<8.
\]

\Section{5. Параллельные прямые\\и проверка утверждений}

При пересечении двух параллельных прямых секущей:

\begin{itemize}
  \item соответственные углы равны;
  \item накрест лежащие углы равны;
  \item сумма внутренних односторонних углов равна $180\degree$.
\end{itemize}

\Key{Обратные признаки}
позволяют установить параллельность по равенству соответственных или
накрест лежащих углов, а также по сумме $180\degree$
внутренних односторонних углов.

\Subsection{Как проверять слова «всегда», «любой», «существует»}

\begin{itemize}
  \item
  Для утверждения со словом «всегда» достаточно одного корректного
  контрпримера, чтобы признать утверждение ложным.

  \item
  Для утверждения «существует» достаточно построить один подходящий
  пример и проверить все условия.

  \item
  Чертёж помогает найти идею.\\
  Окончательный вывод должен опираться на свойства и вычисления.
\end{itemize}

\Key{Пример.}
Равнобедренный треугольник, один угол которого вдвое больше другого,
существует. При углах $x$, $2x$, $2x$ имеем
\[
  5x=180\degree,
  \qquad
  x=36\degree,
\]
то есть получаем допустимые углы
\[
  36\degree,\qquad72\degree,\qquad72\degree.
\]

\clearpage

% ---------- Отдельная страница алгоритма ----------
\thispagestyle{plain}

\begin{center}
  {\Large\bfseries\color{darkaccent}
   Алгоритм: как начать геометрическую задачу}\par

  \vspace{10mm}

\begin{tikzpicture}[node distance=8mm]
  \node[flowstart] (start)
    {Прочитать условие};

  \node[flowstep,below=of start] (mark)
    {Перевести текст на язык чертежа:\\
     отметить равенства, углы, биссектрисы, медианы и перпендикуляры};

  \node[flowchoice,below=10mm of mark] (target)
    {Что требуется найти?};

  \node[
    flowstep,
    below left=12mm and -5mm of target,
    text width=55mm
  ] (angle)
    {Угол: искать треугольник, смежную пару или связи
     при параллельных прямых};

  \node[
    flowstep,
    below right=12mm and -5mm of target,
    text width=55mm
  ] (side)
    {Сторона: искать равные треугольники
     или прямоугольный треугольник};

  \node[flowchoice,below=25mm of target] (enough)
    {Данных достаточно?};

  \node[flowstep,below=11mm of enough] (aux)
    {Если данных мало, выполнить полезное построение
     или выразить углы через $x$};

  \node[flowstep,below=of aux] (solve)
    {Записать равенство или уравнение, вычислить
     и назвать использованное свойство};

  \node[flowstart,below=of solve] (check)
    {Проверить результат};

  \draw[flowarrow] (start) -- (mark);
  \draw[flowarrow] (mark) -- (target);

  \draw[flowarrow]
    (target) -|
    node[pos=0.25,above,font=\small]{угол}
    (angle);

  \draw[flowarrow]
    (target) -|
    node[pos=0.25,above,font=\small]{сторона}
    (side);

  \draw[flowarrow] (angle) |- (enough);
  \draw[flowarrow] (side) |- (enough);

  \draw[flowarrow]
    (enough) --
    node[right,font=\small]{нет}
    (aux);

  \draw[flowarrow] (aux) -- (solve);
  \draw[flowarrow] (solve) -- (check);

  \draw[flowarrow]
    (enough.east) --
    ++(22mm,0)
    node[above,font=\small]{да}
    |- (solve.east);
\end{tikzpicture}
\end{center}

\clearpage

\Section{6. Статистика}

Пусть значения
\[
  x_1,x_2,\ldots,x_n
\]
образуют выборку.

\Subsection{Среднее арифметическое}

\[
  \overline{x}
  =
  \frac{x_1+x_2+\cdots+x_n}{n}.
\]

Если данные представлены таблицей частот, используйте произведения
«значение $\times$ частота»:
\[
  \overline{x}
  =
  \frac{x_1n_1+x_2n_2+\cdots+x_kn_k}
       {n_1+n_2+\cdots+n_k}.
\]

\Subsection{Медиана и мода}

\Key{Медиана} — центральное значение упорядоченной выборки.\\
При чётном числе элементов берут среднее двух центральных значений.

\Key{Мода} — значение, которое встречается чаще остальных.
Мод может быть несколько; иногда мода отсутствует.

\Key{Пример.}
Оценки $2$, $3$, $4$, $5$ получены соответственно
$1$, $4$, $7$, $3$ раза.

Объём выборки равен
\[
  1+4+7+3=15.
\]

Мода равна $4$, а среднее:
\[
  \overline{x}
  =
  \frac{2\cdot1+3\cdot4+4\cdot7+5\cdot3}{15}
  =
  \frac{57}{15}
  =
  3{,}8.
\]

Восьмой элемент упорядоченной выборки равен $4$, поэтому медиана равна $4$.

\Subsection{Круговая диаграмма}

Если величина категории равна $a$, а сумма всех категорий равна $S$, то её
доля и центральный угол сектора вычисляются по формулам
\[
  p=\frac{a}{S},
  \qquad
  p_{\%}=\frac{a}{S}\cdot100\%,
  \qquad
  \varphi=\frac{a}{S}\cdot360\degree.
\]

Для выбора диаграммы сначала сравните крупнейшую и наименьшую доли.
Затем проверьте близкие варианты ещё по одному сектору.

\Key{Быстрая проверка сектора.}
Если $a=31$, $S=160$, то
\[
  p_{\%}
  =
  \frac{31}{160}\cdot100\%
  \approx19{,}4\%,
  \qquad
  \varphi
  =
  \frac{31}{160}\cdot360\degree
  =
  69{,}75\degree.
\]

Процентная доля и угол сектора должны описывать одну и ту же часть целого.

\clearpage

\Section{7. Классическая вероятность}

Если все исходы равновозможны, то
\[
  P(A)=\frac{m}{n},
\]
где $n$ — число всех исходов, а $m$ — число исходов,
благоприятных событию $A$.

Для противоположного события:
\[
  P(\overline{A})=1-P(A).
\]

\Subsection{Два броска симметричной монеты}

\begin{center}
\begin{tabular}{@{}c cc@{}}
\toprule
 & \textbf{Орёл} & \textbf{Решка} \\
\midrule
\textbf{Орёл} & ОО & ОР \\
\textbf{Решка} & РО & РР \\
\bottomrule
\end{tabular}
\end{center}

Всего четыре равновозможных исхода.
Ровно один орёл появляется в исходах ОР и РО, поэтому
\[
  P(\text{ровно один орёл})
  =
  \frac{2}{4}
  =
  0{,}5.
\]

\Key{Пример с противоположным событием.}
Вероятность получить хотя бы один орёл удобнее найти через единственный
противоположный исход РР:
\[
  P(\text{хотя бы один орёл})
  =
  1-P(\text{РР})
  =
  1-\frac14
  =
  \frac34.
\]

\Key{Типичные ошибки:}

\begin{itemize}
  \item перепутать число благоприятных исходов с числом всех исходов;
  \item пропустить порядок результатов: ОР и РО — разные исходы;
  \item записать процент вместо вероятности, когда требуется десятичная дробь;
  \item считать исходы равновозможными без проверки условия.
\end{itemize}

\clearpage

\Section{8. Разобранные задачи с занятия}

\Subsection{Задача 1. Внутренние и внешние углы}

\Key{Условие.}
В треугольнике $ABC$ угол $B$ равен $62\degree$.
Внешний угол при вершине $A$ равен $138\degree$.
Найдите внешний угол при вершине $C$.

\Key{Решение.}
Внутренний угол при вершине $A$ смежный с внешним:
\[
  \angle BAC
  =
  180\degree-138\degree
  =
  42\degree.
\]

Тогда
\[
  \angle BCA
  =
  180\degree-62\degree-42\degree
  =
  76\degree.
\]

Искомый внешний угол смежный с $\angle BCA$:
\[
  180\degree-76\degree
  =
  104\degree.
\]

Тот же результат даёт теорема о внешнем угле:
\[
  42\degree+62\degree=104\degree.
\]

\Key{Ответ:} $104\degree$.

\Subsection{Задача 2. Два внешних угла}

\Key{Условие.}
Внешние углы треугольника $ABC$ при вершинах $A$ и $C$
равны $154\degree$ и $50\degree$ соответственно.
Найдите $\angle ABC$.

\Key{Решение.}
\[
  \angle BAC
  =
  180\degree-154\degree
  =
  26\degree,
\]
\[
  \angle BCA
  =
  180\degree-50\degree
  =
  130\degree.
\]

Следовательно,
\[
  \angle ABC
  =
  180\degree-26\degree-130\degree
  =
  24\degree.
\]

\Key{Ответ:} $24\degree$.

\clearpage

\Section{Разобранные задачи:\\равнобедренный треугольник}

\Subsection{Задача 3. Биссектриса угла при основании}

\Key{Условие.}
В равнобедренном треугольнике $ABC$ с основанием $AC$
проведена биссектриса $AD$, где $D$ лежит на стороне $BC$.
Известно, что
\[
  \angle ADC=132\degree.
\]
Найдите $\angle CBA$.

\Key{Решение.}
Пусть
\[
  \angle BAD=\angle DAC=x.
\]

Тогда
\[
  \angle BAC=2x.
\]

Углы при основании равны. Следовательно,
\[
  \angle BCA=2x.
\]

В треугольнике $ADC$:
\[
  x+2x+132\degree=180\degree,
  \qquad
  x=16\degree.
\]

Значит,
\[
  \angle BAC=\angle BCA=32\degree,
\]
поэтому
\[
  \angle CBA
  =
  180\degree-32\degree-32\degree
  =
  116\degree.
\]

\Key{Ответ:} $116\degree$.

\Subsection{Задача 4. Биссектриса и параллельные прямые}

\Key{Условие.}
В треугольнике $CDE$ точка $F$ лежит на стороне $DE$,
а $CF$ — биссектриса угла $DCE$.
Через вершину $C$ проведена прямая
\[
  AB\parallel DE.
\]

Известно, что
\[
  \angle CDF=54\degree,
  \qquad
  \angle CEF=62\degree.
\]

Найдите $\angle ACF$.

\begin{center}
\begin{tikzpicture}[scale=0.68]
  \tkzDefPoint(0,0){D}
  \tkzDefPoint(6,0){E}
  \tkzDefPoint(3.465,4.767){C}
  \tkzDefPoint(3.13,0){F}
  \tkzDefPoint(-0.5,4.767){A}
  \tkzDefPoint(7.0,4.767){B}

  \tkzDrawSegments[thick,color=accent](D,E D,C C,E C,F A,B)
  \tkzDrawPoints(D,E,C,F,A,B)

  \tkzLabelPoints[below](D,F,E)
  \tkzLabelPoints[above](A,C,B)

  \tkzMarkAngle[size=0.65,color=darkaccent](D,C,F)
  \tkzMarkAngle[size=0.78,color=darkaccent](F,C,E)
\end{tikzpicture}
\end{center}

\Key{Решение.}
Так как $F\in DE$, имеем
\[
  \angle CDE=54\degree,
  \qquad
  \angle CED=62\degree.
\]

Поэтому
\[
  \angle DCE
  =
  180\degree-54\degree-62\degree
  =
  64\degree.
\]

Биссектриса делит его пополам:
\[
  \angle DCF
  =
  \angle FCE
  =
  32\degree.
\]

Из параллельности $AB\parallel DE$ следует
\[
  \angle ACD=\angle CDE=54\degree.
\]

Луч $CD$ лежит внутри угла $ACF$, поэтому
\[
  \angle ACF
  =
  \angle ACD+\angle DCF
  =
  54\degree+32\degree
  =
  86\degree.
\]

\Key{Ответ:} $86\degree$.

\Key{Самопроверка.}
\[
  54\degree+62\degree+64\degree=180\degree,
\]
а
\[
  54\degree+\frac{64\degree}{2}=86\degree.
\]

\Section{Разобранные задачи: продолжение стороны}

\Subsection{Задача 5. Равные стороны в двух треугольниках}

\Key{Условие.}
В равнобедренном треугольнике $ABC$ основанием является $AC$.
Точка $D$ лежит на продолжении стороны $AB$ за вершину $B$, причём
\[
  BD=BC.
\]

Известно, что
\[
  \angle ADC=55\degree.
\]

Найдите $\angle ABC$.

\begin{center}
\begin{tikzpicture}[scale=0.92]
  \tkzDefPoint(-4,0){A}
  \tkzDefPoint(0,0){B}
  \tkzDefPoint(4,0){D}
  \tkzDefPoint(1.368,-3.759){C}

  \tkzDrawSegments[thick,color=accent](A,D A,C B,C C,D)
  \tkzDrawPoints(A,B,C,D)

  \tkzLabelPoints[above](A,B,D)
  \tkzLabelPoints[below](C)

  \tkzMarkSegments[mark=|,color=darkaccent](A,B B,C B,D)
  \tkzMarkAngle[size=0.62,color=darkaccent](B,D,C)
  \tkzLabelAngle[pos=0.92](B,D,C){$55\degree$}
\end{tikzpicture}
\end{center}

\Key{Решение.}
Лучи $DA$ и $DB$ совпадают, поэтому
\[
  \angle BDC=\angle ADC=55\degree.
\]

В треугольнике $BDC$ стороны $BD$ и $BC$ равны,
значит углы при основании $DC$ равны:
\[
  \angle BDC=\angle BCD=55\degree.
\]

Отсюда
\[
  \angle DBC
  =
  180\degree-55\degree-55\degree
  =
  70\degree.
\]

Лучи $BA$ и $BD$ противоположны, следовательно,
\[
  \angle ABC+\angle DBC=180\degree.
\]

Значит,
\[
  \angle ABC
  =
  180\degree-70\degree
  =
  110\degree.
\]

\Key{Ответ:} $110\degree$.

\Needspace{15\baselineskip}

\Subsection{Задача 6. Две связанные модели с углом $30\degree$}

\Key{Условие.}
В треугольнике $ABC$ выполнено
\[
  AB=BC,
  \qquad
  \angle ABC=120\degree.
\]

Отрезок $BM$ — медиана.
На луче $BM$ за точкой $M$ выбрана точка $F$ так, что
\[
  \angle BAF=90\degree.
\]

Известно, что
\[
  FM=63.
\]

Найдите $AB$.

\begin{center}
\begin{tikzpicture}[scale=0.72]
  \tkzDefPoint(-3.464,0){A}
  \tkzDefPoint(3.464,0){C}
  \tkzDefPoint(0,2){B}
  \tkzDefPoint(0,0){M}
  \tkzDefPoint(0,-6){F}

  \tkzDrawSegments[thick,color=accent](A,B B,C A,C B,F A,F)
  \tkzDrawPoints(A,B,C,M,F)

  \tkzLabelPoints[left](A)
  \tkzLabelPoints[right](C)
  \tkzLabelPoints[above](B)
  \tkzLabelPoints[right](M,F)

  \tkzMarkSegments[mark=|,color=darkaccent](A,B B,C)
  \tkzMarkSegments[mark=||,color=darkaccent](A,M M,C)

  \tkzMarkRightAngle[size=0.28,color=darkaccent](B,M,A)
  \tkzMarkRightAngle[size=0.28,color=darkaccent](B,A,F)
\end{tikzpicture}
\end{center}

\Key{Решение.}
Пусть
\[
  AB=x.
\]

В равнобедренном треугольнике медиана $BM$ является биссектрисой и высотой.
Поэтому
\[
  \angle ABM=60\degree,
  \qquad
  \angle BAM=30\degree.
\]

В прямоугольном треугольнике $ABM$ катет $BM$ лежит напротив угла
$30\degree$, следовательно,
\[
  BM=\frac{x}{2}.
\]

Точки $B$, $M$, $F$ лежат на одной прямой, поэтому
\[
  \angle ABF=60\degree.
\]

В прямоугольном треугольнике $ABF$ получаем
\[
  \angle AFB=30\degree.
\]

Значит, катет $AB$ равен половине гипотенузы $BF$:
\[
  x=\frac{BF}{2}.
\]

При этом
\[
  BF
  =
  BM+MF
  =
  \frac{x}{2}+63.
\]

Получаем уравнение
\[
  x
  =
  \frac{\frac{x}{2}+63}{2}.
\]

Следовательно,
\[
  2x=\frac{x}{2}+63,
\]
\[
  \frac{3x}{2}=63,
\]
\[
  x=42.
\]

\Key{Ответ:} $42$.

\clearpage

\Section{Разобранные задачи: проверка утверждений}

В заданиях такого типа каждое утверждение рассматривается отдельно.
Слова «любой» и «всегда» требуют особенно строгой проверки.

\begin{tabularx}{\linewidth}{
  @{}
  >{\bfseries\color{darkaccent}\raggedright\arraybackslash}
  p{0.30\linewidth}
  >{\raggedright\arraybackslash}
  X
  @{}
}
\toprule
Опорный факт & Как его использовать \\
\midrule

Равнобедренный треугольник &
Равным сторонам соответствуют равные углы, а равным углам — равные стороны. \\

Треугольник с углами
$36\degree$, $72\degree$, $72\degree$ &
Даёт пример существования равнобедренного треугольника,
один угол которого вдвое больше другого. \\

Неравенство треугольника &
Каждая сторона меньше суммы двух других сторон. \\

Прямоугольный треугольник с углом $30\degree$ &
Катет напротив $30\degree$ равен половине гипотенузы. \\

Перпендикуляры к одной прямой &
Две прямые, перпендикулярные одной и той же прямой, параллельны. \\

Касательная к окружности &
Радиус, проведённый в точку касания, перпендикулярен касательной. \\

Биссектриса угла &
Делит угол пополам; точка на биссектрисе равноудалена от сторон угла. \\

Признак параллельности &
Равенство накрест лежащих или соответственных углов
позволяет доказать параллельность. \\

\bottomrule
\end{tabularx}

\Subsection{Отдельные ловушки}

\begin{itemize}
  \item
  Вертикальные углы равны, однако их сумма обычно отличается от
  $180\degree$.

  \item
  В общем треугольнике биссектриса угла не обязана быть медианой.

  \item
  Через точку вне окружности можно провести ровно две касательные
  к окружности.

  \item
  Теорема Пифагора содержит сумму квадратов катетов:
  \[
    c^2=a^2+b^2.
  \]

  \item
  В любом треугольнике хотя бы один угол не превосходит $60\degree$:
  иначе сумма трёх углов была бы больше $180\degree$.
\end{itemize}

\Section{Разобранные задачи: статистика}
\enlargethispage{2\baselineskip}

\Subsection{Задача 7. Восстановление значения по среднему}

\Key{Условие.}
Пять объёмов равны
\[
  23,\quad24,\quad25,\quad x,\quad57,
\]
а их среднее арифметическое равно $32$.
Найдите $x$ и оцените процентные доли всех пяти значений.

\Key{Решение.}
Общая сумма значений:
\[
  32\cdot5=160.
\]

Сумма четырёх известных значений:
\[
  23+24+25+57=129.
\]

Поэтому
\[
  x=160-129=31.
\]

Процентные доли:
\[
\begin{aligned}
  \frac{23}{160}\cdot100\%
  &\approx14{,}4\%,
  &
  \frac{24}{160}\cdot100\%
  &=15\%,
  \\
  \frac{25}{160}\cdot100\%
  &\approx15{,}6\%,
  &
  \frac{31}{160}\cdot100\%
  &\approx19{,}4\%,
  \\
  \frac{57}{160}\cdot100\%
  &\approx35{,}6\%.
\end{aligned}
\]

На круговой диаграмме первые три сектора должны быть близки по размеру,
а крупнейший сектор должен занимать немного больше трети круга.

\Key{Ответ:}
\[
  x=31;
\]
доли приблизительно
\[
  14{,}4\%,\quad
  15\%,\quad
  15{,}6\%,\quad
  19{,}4\%,\quad
  35{,}6\%.
\]


\newpage
\Subsection{Задача 8. Таблица частот}

\Key{Условие.}
Оценки $2$, $3$, $4$, $5$ получили соответственно
$1$, $4$, $7$, $3$ ученика.
Найдите среднюю оценку, медиану и моду.

\Key{Решение.}
Объём выборки:
\[
  n=1+4+7+3=15.
\]

Среднее:
\[
  \overline{x}
  =
  \frac{2\cdot1+3\cdot4+4\cdot7+5\cdot3}{15}
  =
  \frac{57}{15}
  =
  3{,}8.
\]

При $15$ значениях медиана занимает восьмое место.
На местах с шестого по двенадцатое стоит оценка $4$,
поэтому медиана равна $4$.

Чаще всего также встречается оценка $4$, значит мода равна $4$.

\Key{Ответ:}
среднее $3{,}8$; медиана $4$; мода $4$.

\clearpage

\Section{Разобранные задачи: закрепление статистики и вероятности}

\Subsection{Задача 9. Ещё одна таблица частот}

\Key{Условие.}
Оценки $2$, $3$, $4$, $5$ получили соответственно
$2$, $3$, $6$, $1$ учащийся.
Найдите среднюю оценку, медиану и моду.

\Key{Решение.}
\[
  n=2+3+6+1=12.
\]

Среднее:
\[
  \overline{x}
  =
  \frac{2\cdot2+3\cdot3+4\cdot6+5\cdot1}{12}
  =
  \frac{42}{12}
  =
  3{,}5.
\]

Шестое и седьмое значения равны $4$, поэтому медиана равна $4$.
Наибольшая частота равна $6$, следовательно, мода равна $4$.

\Key{Ответ:}
среднее $3{,}5$; медиана $4$; мода $4$.

\Subsection{Задача 10. Шары в мешке}

\Key{Условие.}
В мешке находятся $5$ белых, $3$ чёрных и $2$ серых шара.
Наугад выбирают один шар.
Найдите вероятность того, что он окажется чёрным,
и вероятность того, что он окажется любого цвета, кроме чёрного.

\Key{Решение.}
Всего шаров
\[
  5+3+2=10.
\]

Поэтому
\[
  P(\text{чёрный})
  =
  \frac{3}{10}
  =
  0{,}3.
\]

Подходящих шаров любого цвета, кроме чёрного, семь:
\[
  P(\text{любой, кроме чёрного})
  =
  \frac{7}{10}
  =
  0{,}7.
\]

Также можно использовать противоположное событие:
\[
  1-0{,}3=0{,}7.
\]

\Key{Ответ:} $0{,}3$ и $0{,}7$.

\clearpage

\Subsection{Задача 11. Два броска монеты}

\Key{Условие.}
Симметричную монету бросают два раза.
Найдите вероятность того, что орёл выпадет ровно один раз.

\Key{Решение.}
Равновозможные исходы:
\[
  \text{ОО},
  \quad
  \text{ОР},
  \quad
  \text{РО},
  \quad
  \text{РР}.
\]

Событию подходят исходы ОР и РО. Поэтому
\[
  P
  =
  \frac{2}{4}
  =
  \frac12
  =
  0{,}5.
\]

\Key{Ответ:} $0{,}5$.

% ---------- Тренировочный блок ----------
\Section{Тренировочный блок — Путь А}

Решите 20 заданий.
Задания расположены по нарастающей сложности.
В блоке геометрии записывайте краткое обоснование каждого существенного шага.

\begin{enumerate}[label=\textbf{\arabic*.}]

  \item
  В треугольнике $ABC$ углы $A$ и $B$ равны
  $47\degree$ и $58\degree$.
  Сторона $AC$ продолжена за точку $C$ до точки $D$.
  Найдите $\angle BCD$.

  \item
  Внешний угол при вершине $A$ треугольника $ABC$ равен $142\degree$,
  а внешний угол при вершине $C$ равен $133\degree$.
  Найдите $\angle ABC$.

  \item
  В равнобедренном треугольнике $ABC$ выполнено $AB=BC$.
  Биссектриса $AD$ пересекает сторону $BC$, причём
  \[
    \angle ADC=126\degree.
  \]
  Найдите $\angle ABC$.

  \item
  В треугольнике $CDE$ отрезок $CF$ — биссектриса угла $DCE$,
  точка $F$ лежит на $DE$.
  Известно, что
  \[
    \angle CDE=48\degree,
    \qquad
    \angle CED=64\degree.
  \]
  Найдите $\angle DFC$.

  \item
  В равнобедренном треугольнике $ABC$ выполнено $AB=BC$.
  Внешний угол при вершине $A$ равен $132\degree$.
  Найдите $\angle ABC$.

  \item
  Две стороны треугольника равны $5$ и $8$.
  Сколько целых положительных значений может принимать третья сторона?

  \item
  В прямоугольном треугольнике катет, лежащий напротив угла $30\degree$,
  равен $7$.
  Найдите гипотенузу.

  \item
  Гипотенуза прямоугольного треугольника равна $18$,
  один из острых углов равен $30\degree$.
  Найдите меньший катет.

  \item
  Один из острых углов прямоугольного треугольника равен $45\degree$,
  а один катет равен $6$.
  Найдите второй катет и второй острый угол.

  \item
  Две параллельные прямые пересечены секущей.
  Один из внутренних односторонних углов равен $124\degree$.
  Найдите второй такой угол.

  \item
  Радиус $OT$ проведён в точку касания $T$ прямой и окружности.
  Найдите угол между $OT$ и касательной.

  \item
  Точка $M$ лежит на биссектрисе угла.
  Расстояние от $M$ до одной стороны угла равно $5{,}4$ см.
  Найдите расстояние от $M$ до другой стороны.

  \item
  В двух треугольниках равны соответственно две стороны и угол между ними.
  По какому признаку равны эти треугольники?

  \item
  В равнобедренном треугольнике $ABC$ выполнено
  \[
    AB=BC,
    \qquad
    \angle ABC=120\degree.
  \]
  Медиана $BM$ проведена к основанию $AC$.
  На луче $BM$ за точкой $M$ выбрана точка $F$ так, что
  \[
    AF\perp AB.
  \]
  Известно, что
  \[
    FM=54.
  \]
  Найдите $AB$.

  \item
  Выберите верные утверждения:
  \begin{enumerate}[label=\arabic*)]
    \item
    существует равнобедренный треугольник, один угол которого
    вдвое больше другого;

    \item
    сумма любой пары вертикальных углов равна $180\degree$;

    \item
    биссектриса любого треугольника делит противоположную сторону пополам;

    \item
    длина каждой стороны треугольника меньше суммы двух других сторон.
  \end{enumerate}

  \item
  Оценки $2$, $3$, $4$, $5$ получили соответственно
  $2$, $3$, $6$, $1$ учащихся.
  Найдите среднюю оценку, медиану и моду.

  \item
  Пять чисел равны
  \[
    21,\quad24,\quad29,\quad x,\quad61,
  \]
  а их среднее арифметическое равно $36$.
  Найдите $x$.

  \item
  Объёмы пяти водохранилищ пропорциональны числам
  \[
    23,\quad24,\quad25,\quad31,\quad57.
  \]
  Найдите процентную долю самого большого и самого маленького значения.
  Ответ округлите до десятых процента.

  \item
  В мешке лежат $4$ чёрных, $6$ белых и $5$ синих шаров.
  Наугад вынимают один шар.
  Найдите вероятность того, что он:
  а) чёрный;
  б) любого цвета, кроме чёрного.

  \item
  Симметричную монету бросают два раза.
  Найдите вероятность того, что орёл выпадет ровно один раз.

\end{enumerate}

% ---------- Ответы ----------
\Section{Ответы}

\footnotesize
\renewcommand{\arraystretch}{1.00}

\begin{longtable}{
  @{}
  >{\bfseries\color{darkaccent}}p{10mm}
  p{0.82\linewidth}
  @{}
}
\toprule
№ & Ответ \\
\midrule
\endfirsthead

\toprule
№ & Ответ \\
\midrule
\endhead

1  & $105\degree$. \\

2  & $95\degree$. \\

3  & $108\degree$. \\

4  & $98\degree$. \\

5  & $84\degree$. \\

6  & $9$ значений:
     $4,5,6,7,8,9,10,11,12$. \\

7  & $14$. \\

8  & $9$. \\

9  & Второй катет $6$;
     второй острый угол $45\degree$. \\

10 & $56\degree$. \\

11 & $90\degree$. \\

12 & $5{,}4$ см. \\

13 & По первому признаку:
     две стороны и угол между ними. \\

14 & $36$. \\

15 & Верны утверждения 1 и 4. \\

16 & Среднее $3{,}5$;
     медиана $4$;
     мода $4$. \\

17 & $45$. \\

18 & Наибольшая доля $35{,}6\%$;
     наименьшая доля $14{,}4\%$. \\

19 & а) $\dfrac{4}{15}$;
     б) $\dfrac{11}{15}$. \\

20 & $\dfrac12=0{,}5$. \\

\bottomrule
\end{longtable}

\normalsize

\end{document}